% prezentace.tex -- Lekce 07: Seznamy a hudba
% preamble-prez.tex -- sdílená preamble pro prezentace (beamer)
% Includuje se z ucitel/lekceXX/prezentace.tex jako:
%   % preamble-prez.tex -- sdílená preamble pro prezentace (beamer)
% Includuje se z ucitel/lekceXX/prezentace.tex jako:
%   % preamble-prez.tex -- sdílená preamble pro prezentace (beamer)
% Includuje se z ucitel/lekceXX/prezentace.tex jako:
%   \input{../spolecne/preamble-prez.tex}

\documentclass[aspectratio=169]{beamer}

% --- Jazyk a font ---
\usepackage{polyglossia}
\setmainlanguage{czech}
\usepackage{fontspec}

% --- Téma ---
\usetheme{default}
\usecolortheme{default}

\definecolor{microbit-blue}{HTML}{1E4D8C}
\definecolor{microbit-lightblue}{HTML}{D6E4F7}
\definecolor{code-bg}{HTML}{F0F0F0}

\setbeamercolor{frametitle}{bg=microbit-blue, fg=white}
\setbeamercolor{title}{fg=microbit-blue}
\setbeamercolor{block title}{bg=microbit-blue, fg=white}
\setbeamercolor{block body}{bg=microbit-lightblue}
\setbeamercolor{itemize item}{fg=microbit-blue}
\setbeamercolor{itemize subitem}{fg=microbit-blue!70}
\setbeamerfont{frametitle}{size=\large, series=\bfseries}
\setbeamertemplate{navigation symbols}{}
\setbeamertemplate{footline}{%
  \leavevmode%
  \hbox{%
    \begin{beamercolorbox}[wd=.5\paperwidth,ht=2.5ex,dp=1.125ex,leftskip=6pt]{author in head/foot}%
      \usebeamerfont{author in head/foot}\textcolor{gray}{\small Úvod do Pythonu na Micro:bitu}
    \end{beamercolorbox}%
    \begin{beamercolorbox}[wd=.5\paperwidth,ht=2.5ex,dp=1.125ex,rightskip=6pt,right]{date in head/foot}%
      \usebeamerfont{date in head/foot}\textcolor{gray}{\small\insertframenumber\,/\,\inserttotalframenumber}
    \end{beamercolorbox}}%
}

% --- Česky korektní uvozovky ---
\usepackage{csquotes}
\newcommand{\uv}[1]{\enquote{#1}}

% --- Zdrojový kód (Python) ---
\usepackage{listings}
\lstset{
  language=Python,
  basicstyle=\ttfamily\small,
  backgroundcolor=\color{code-bg},
  frame=single,
  framerule=0pt,
  xleftmargin=6pt,
  xrightmargin=6pt,
  breaklines=true,
  showstringspaces=false,
  keywordstyle=\color{microbit-blue}\bfseries,
  stringstyle=\color{teal},
  commentstyle=\color{gray}\itshape,
  literate=
    {á}{{\'a}}1 {é}{{\'e}}1 {í}{{\'i}}1 {ó}{{\'o}}1 {ú}{{\'u}}1
    {ů}{{\r{u}}}1 {č}{{\v{c}}}1 {š}{{\v{s}}}1 {ž}{{\v{z}}}1
    {Á}{{\'A}}1 {É}{{\'E}}1 {Í}{{\'I}}1 {Ó}{{\'O}}1 {Ú}{{\'U}}1
    {Č}{{\v{C}}}1 {Š}{{\v{S}}}1 {Ž}{{\v{Z}}}1 {ň}{{\v{n}}}1 {ř}{{\v{r}}}1,
}

% --- Zvýraznění pojmů ---
\newcommand{\pojem}[1]{\textbf{\textcolor{microbit-blue}{#1}}}
\newcommand{\pyinline}[1]{\texttt{#1}}

% --- Grafika a diagramy ---
\usepackage{tikz}

% --- Ostatní ---
\usepackage{graphicx}
\usepackage{hyperref}
\hypersetup{colorlinks=true, urlcolor=microbit-blue}


\documentclass[aspectratio=169]{beamer}

% --- Jazyk a font ---
\usepackage{polyglossia}
\setmainlanguage{czech}
\usepackage{fontspec}

% --- Téma ---
\usetheme{default}
\usecolortheme{default}

\definecolor{microbit-blue}{HTML}{1E4D8C}
\definecolor{microbit-lightblue}{HTML}{D6E4F7}
\definecolor{code-bg}{HTML}{F0F0F0}

\setbeamercolor{frametitle}{bg=microbit-blue, fg=white}
\setbeamercolor{title}{fg=microbit-blue}
\setbeamercolor{block title}{bg=microbit-blue, fg=white}
\setbeamercolor{block body}{bg=microbit-lightblue}
\setbeamercolor{itemize item}{fg=microbit-blue}
\setbeamercolor{itemize subitem}{fg=microbit-blue!70}
\setbeamerfont{frametitle}{size=\large, series=\bfseries}
\setbeamertemplate{navigation symbols}{}
\setbeamertemplate{footline}{%
  \leavevmode%
  \hbox{%
    \begin{beamercolorbox}[wd=.5\paperwidth,ht=2.5ex,dp=1.125ex,leftskip=6pt]{author in head/foot}%
      \usebeamerfont{author in head/foot}\textcolor{gray}{\small Úvod do Pythonu na Micro:bitu}
    \end{beamercolorbox}%
    \begin{beamercolorbox}[wd=.5\paperwidth,ht=2.5ex,dp=1.125ex,rightskip=6pt,right]{date in head/foot}%
      \usebeamerfont{date in head/foot}\textcolor{gray}{\small\insertframenumber\,/\,\inserttotalframenumber}
    \end{beamercolorbox}}%
}

% --- Česky korektní uvozovky ---
\usepackage{csquotes}
\newcommand{\uv}[1]{\enquote{#1}}

% --- Zdrojový kód (Python) ---
\usepackage{listings}
\lstset{
  language=Python,
  basicstyle=\ttfamily\small,
  backgroundcolor=\color{code-bg},
  frame=single,
  framerule=0pt,
  xleftmargin=6pt,
  xrightmargin=6pt,
  breaklines=true,
  showstringspaces=false,
  keywordstyle=\color{microbit-blue}\bfseries,
  stringstyle=\color{teal},
  commentstyle=\color{gray}\itshape,
  literate=
    {á}{{\'a}}1 {é}{{\'e}}1 {í}{{\'i}}1 {ó}{{\'o}}1 {ú}{{\'u}}1
    {ů}{{\r{u}}}1 {č}{{\v{c}}}1 {š}{{\v{s}}}1 {ž}{{\v{z}}}1
    {Á}{{\'A}}1 {É}{{\'E}}1 {Í}{{\'I}}1 {Ó}{{\'O}}1 {Ú}{{\'U}}1
    {Č}{{\v{C}}}1 {Š}{{\v{S}}}1 {Ž}{{\v{Z}}}1 {ň}{{\v{n}}}1 {ř}{{\v{r}}}1,
}

% --- Zvýraznění pojmů ---
\newcommand{\pojem}[1]{\textbf{\textcolor{microbit-blue}{#1}}}
\newcommand{\pyinline}[1]{\texttt{#1}}

% --- Grafika a diagramy ---
\usepackage{tikz}

% --- Ostatní ---
\usepackage{graphicx}
\usepackage{hyperref}
\hypersetup{colorlinks=true, urlcolor=microbit-blue}


\documentclass[aspectratio=169]{beamer}

% --- Jazyk a font ---
\usepackage{polyglossia}
\setmainlanguage{czech}
\usepackage{fontspec}

% --- Téma ---
\usetheme{default}
\usecolortheme{default}

\definecolor{microbit-blue}{HTML}{1E4D8C}
\definecolor{microbit-lightblue}{HTML}{D6E4F7}
\definecolor{code-bg}{HTML}{F0F0F0}

\setbeamercolor{frametitle}{bg=microbit-blue, fg=white}
\setbeamercolor{title}{fg=microbit-blue}
\setbeamercolor{block title}{bg=microbit-blue, fg=white}
\setbeamercolor{block body}{bg=microbit-lightblue}
\setbeamercolor{itemize item}{fg=microbit-blue}
\setbeamercolor{itemize subitem}{fg=microbit-blue!70}
\setbeamerfont{frametitle}{size=\large, series=\bfseries}
\setbeamertemplate{navigation symbols}{}
\setbeamertemplate{footline}{%
  \leavevmode%
  \hbox{%
    \begin{beamercolorbox}[wd=.5\paperwidth,ht=2.5ex,dp=1.125ex,leftskip=6pt]{author in head/foot}%
      \usebeamerfont{author in head/foot}\textcolor{gray}{\small Úvod do Pythonu na Micro:bitu}
    \end{beamercolorbox}%
    \begin{beamercolorbox}[wd=.5\paperwidth,ht=2.5ex,dp=1.125ex,rightskip=6pt,right]{date in head/foot}%
      \usebeamerfont{date in head/foot}\textcolor{gray}{\small\insertframenumber\,/\,\inserttotalframenumber}
    \end{beamercolorbox}}%
}

% --- Česky korektní uvozovky ---
\usepackage{csquotes}
\newcommand{\uv}[1]{\enquote{#1}}

% --- Zdrojový kód (Python) ---
\usepackage{listings}
\lstset{
  language=Python,
  basicstyle=\ttfamily\small,
  backgroundcolor=\color{code-bg},
  frame=single,
  framerule=0pt,
  xleftmargin=6pt,
  xrightmargin=6pt,
  breaklines=true,
  showstringspaces=false,
  keywordstyle=\color{microbit-blue}\bfseries,
  stringstyle=\color{teal},
  commentstyle=\color{gray}\itshape,
  literate=
    {á}{{\'a}}1 {é}{{\'e}}1 {í}{{\'i}}1 {ó}{{\'o}}1 {ú}{{\'u}}1
    {ů}{{\r{u}}}1 {č}{{\v{c}}}1 {š}{{\v{s}}}1 {ž}{{\v{z}}}1
    {Á}{{\'A}}1 {É}{{\'E}}1 {Í}{{\'I}}1 {Ó}{{\'O}}1 {Ú}{{\'U}}1
    {Č}{{\v{C}}}1 {Š}{{\v{S}}}1 {Ž}{{\v{Z}}}1 {ň}{{\v{n}}}1 {ř}{{\v{r}}}1,
}

% --- Zvýraznění pojmů ---
\newcommand{\pojem}[1]{\textbf{\textcolor{microbit-blue}{#1}}}
\newcommand{\pyinline}[1]{\texttt{#1}}

% --- Grafika a diagramy ---
\usepackage{tikz}

% --- Ostatní ---
\usepackage{graphicx}
\usepackage{hyperref}
\hypersetup{colorlinks=true, urlcolor=microbit-blue}


\title{Lekce 07: Seznamy a hudba}
\subtitle{Úvod do Pythonu na Micro:bitu}
\date{}

\begin{document}

\begin{frame}
  \titlepage
\end{frame}

% ------------------------------------------------------------------
\begin{frame}[fragile]{Co když potřebuji víc hodnot najednou?}

  \begin{columns}[T]
    \begin{column}{0.44\textwidth}
      \begin{block}{Dosud — jedna hodnota}
        \begin{lstlisting}
jmeno   = "Alice"
cislo   = 42
obrazek = Image.HEART
        \end{lstlisting}
        \small Každá proměnná drží\\
        \textbf{jednu} hodnotu.
      \end{block}
    \end{column}
    \begin{column}{0.52\textwidth}
      \begin{exampleblock}{Nově — kolekce hodnot}
        \begin{lstlisting}
jmena = ["Alice","Bob","Carla"]
obrazky = [Image.HAPPY,
           Image.SAD,
           Image.SURPRISED]
        \end{lstlisting}
        \small Jedna proměnná drží\\
        \textbf{více hodnot najednou.}
      \end{exampleblock}
    \end{column}
  \end{columns}

  \pause\bigskip
  \begin{center}
  \begin{tikzpicture}
    \foreach \i/\val in {0/Alice, 1/Bob, 2/Carla} {
      \draw[thick] (\i*2.7, 0) rectangle (\i*2.7+2.5, 0.75);
      \node at (\i*2.7+1.25, 0.38) {\texttt{\val}};
      \node[font=\footnotesize, text=microbit-blue] at (\i*2.7+1.25, -0.32)
            {\texttt{jmena[\i]}};
    }
    \node[font=\small, left=4pt] at (0, 0.38) {\texttt{jmena}};
  \end{tikzpicture}
  \end{center}

  \pause\medskip
  K~jednotlivým prvkům přistupujeme přes \pojem{index} — číslo pozice.
  \textbf{Indexy začínají od~0!}
\end{frame}

% ------------------------------------------------------------------
\begin{frame}[fragile]{Seznam v Pythonu — syntaxe}

  Prvky mohou být jakéhokoliv typu — řetězce, čísla, obrázky\ldots

  \bigskip
  \begin{lstlisting}
barvy   = ["cervena", "zelena", "modra"]
cisla   = [1, 2, 3, 4, 5]
obrazky = [Image.HAPPY, Image.SAD, Image.SURPRISED]
  \end{lstlisting}

  \bigskip
  \begin{columns}[T]
    \begin{column}{0.46\textwidth}
      \begin{block}{Jak se vytvoří}
        Hranaté závorky \pyinline{[]}\\
        Prvky oddělené čárkou\\
        Může být prázdný: \pyinline{[]}
      \end{block}
    \end{column}
    \begin{column}{0.50\textwidth}
      \begin{block}{Už jsme ho viděli}
        Lekce 4: \pyinline{[Image.HAPPY, \ldots]}\\
        Lekce 5: \pyinline{random.choice(seznam)}\\
        Teď ho konečně \textbf{poznáme pořádně}.
      \end{block}
    \end{column}
  \end{columns}
\end{frame}

% ------------------------------------------------------------------
\begin{frame}[fragile]{Indexování a délka}

  \begin{lstlisting}
barvy = ["cervena", "zelena", "modra"]

barvy[0]    # "cervena"  -- prvni prvek
barvy[1]    # "zelena"
barvy[-1]   # "modra"    -- posledni prvek

len(barvy)  # 3          -- pocet prvku
  \end{lstlisting}

  \bigskip
  \begin{columns}[T]
    \begin{column}{0.46\textwidth}
      \begin{block}{Kladné indexy}
        \pyinline{[0]} = první\\
        \pyinline{[1]} = druhý\\
        \pyinline{[n-1]} = poslední
      \end{block}
    \end{column}
    \begin{column}{0.46\textwidth}
      \begin{block}{Záporné indexy}
        \pyinline{[-1]} = poslední\\
        \pyinline{[-2]} = předposlední\\
        Počítáme od konce
      \end{block}
    \end{column}
  \end{columns}

  \pause\bigskip
  \begin{alertblock}{IndexError}
    Pokud index neexistuje (např. \pyinline{barvy[5]}), Python ohlásí chybu.
  \end{alertblock}
\end{frame}

% ------------------------------------------------------------------
\begin{frame}[fragile]{Procházení seznamu — \texttt{for ... in}}

  Přesně jako v~lekci~4 — ale místo \pyinline{range()} použijeme seznam přímo:

  \bigskip
  \begin{lstlisting}
melodie = ["C4:4", "E4:4", "G4:4"]

for nota in melodie:
    print(nota)      # vypise kazdou notu
  \end{lstlisting}

  \bigskip
  \begin{lstlisting}
obrazky = [Image.HAPPY, Image.SAD, Image.SURPRISED]

for obrazek in obrazky:
    display.show(obrazek)
    sleep(500)
  \end{lstlisting}

  \bigskip
  Proměnná \pyinline{nota} (nebo \pyinline{obrazek}) dostává vždy \textbf{aktuální prvek} —
  smyčka projde celý seznam od začátku do konce.
\end{frame}

% ------------------------------------------------------------------
\begin{frame}[fragile]{Hudba — \texttt{import music}}

  Modul \pyinline{music} umožňuje přehrávat zvuky a melodie:

  \bigskip
  \begin{lstlisting}
import music

music.play(music.NYAN)      # predefinovana melodie
music.play(music.BIRTHDAY)  # dalsi predefinovana
  \end{lstlisting}

  \bigskip
  \begin{block}{Předdefinované melodie (výběr)}
    \pyinline{NYAN}, \pyinline{BIRTHDAY}, \pyinline{DADADADUM},\\
    \pyinline{ENTERTAINER}, \pyinline{FUNK}, \pyinline{BLUES},\\
    \pyinline{JUMP_UP}, \pyinline{JUMP_DOWN}, \pyinline{POWER_UP}
  \end{block}

  \bigskip
  \begin{alertblock}{Zvuk na micro:bitu}
    Zvuk vychází z~GPIO pinu~0 — zapoj krokodýlky ke~sluchátkům nebo reproduktoru.
    Micro:bit~v2 má vestavěný reproduktor!
  \end{alertblock}
\end{frame}

% ------------------------------------------------------------------
\begin{frame}[fragile]{Vlastní melodie — notový zápis}

  Každá nota je řetězec ve formátu \pyinline{"NOTA[OKTAVA]:DELKA"}:

  \bigskip
  \begin{center}
  \begin{tabular}{ll}
    \toprule
    \textbf{Řetězec} & \textbf{Význam} \\
    \midrule
    \pyinline{"C4:4"}  & nota C, oktáva 4, čtvrťová (4 tikety) \\
    \pyinline{"E4:2"}  & nota E, oktáva 4, půlová (2 tikety) \\
    \pyinline{"R:4"}   & pomlka, čtvrťová \\
    \bottomrule
  \end{tabular}
  \end{center}

  \bigskip
  \begin{lstlisting}
import music

melodie = ["C4:4", "D4:4", "E4:4", "F4:4",
           "G4:2", "R:2",  "G4:4"]

music.play(melodie)
  \end{lstlisting}

  \bigskip
  Nota \pyinline{"C"} bez oktávy = posledně použitá oktáva.
  Délky: 1~=~celá, 2~=~půlová, 4~=~čtvrťová, 8~=~osminová.
\end{frame}

% ------------------------------------------------------------------
\begin{frame}[fragile]{Jukebox — výběr melodie tlačítky}

  \begin{lstlisting}
from microbit import *
import music

while True:
    if button_a.was_pressed():
        music.play(music.NYAN)
    elif button_b.was_pressed():
        music.play(music.FUNK)
    sleep(50)
  \end{lstlisting}

  \bigskip
  \begin{exampleblock}{Rozšíření — vlastní melodie v seznamu}
    \begin{lstlisting}
melodie = ["C4:4", "E4:4", "G4:4", "C5:2"]

if button_a.was_pressed():
    music.play(melodie)
    \end{lstlisting}
  \end{exampleblock}
\end{frame}

% ------------------------------------------------------------------
\begin{frame}{Co jsme se naučili}
  \begin{itemize}
    \item[\checkmark] \pojem{Seznam} = proměnná s více hodnotami; přístup přes \pojem{index}
    \item[\checkmark] \pyinline{seznam[0]} — první prvek; \pyinline{seznam[-1]} — poslední
    \item[\checkmark] \pyinline{len(seznam)} — počet prvků
    \item[\checkmark] \pyinline{for prvek in seznam:} — procházení od začátku do konce
    \item[\checkmark] \pyinline{import music} — modul pro přehrávání zvuků
    \item[\checkmark] \pyinline{music.play(melodie)} — přehraje seznam not nebo předdefinovanou melodii
    \item[\checkmark] Notový formát \pyinline{"C4:4"} — nota, oktáva, délka
    \item[\checkmark] Melodie je seznam řetězců — \textbf{seznam a hudba jsou jedno}
  \end{itemize}
\end{frame}

\end{document}
